% AUTHOR: Amlal El Mahrouss
% PURPOSE: WG01: A Methodology for Freestanding Development

\documentclass[11pt, a4paper]{article}
\usepackage{graphicx}
\usepackage{listings}
\usepackage{xcolor}
\usepackage{hyperref}
\usepackage[margin=0.5in,top=1in,bottom=1in]{geometry}

\setlength{\columnsep}{20pt}

\definecolor{codegray}{gray}{0.95}
\definecolor{codeblue}{rgb}{0.1,0.1,0.8}
\definecolor{codegreen}{rgb}{0,0.6,0}
\definecolor{codepurple}{rgb}{0.58,0,0.82}

\title{Methodology for Freestanding Development.}
\author{Amlal El Mahrouss\\amlal@nekernel.org}
\date{December 2025\\Last Edited: January 2026}

\lstset{
    language=C++,
    backgroundcolor=\color{codegray},
    basicstyle=\footnotesize\ttfamily,
    keywordstyle=\color{codeblue}\bfseries,
    commentstyle=\color{codegreen},
    stringstyle=\color{codepurple},
    numbers=left,
    numberstyle=\tiny\color{gray},
    stepnumber=1,
    numbersep=5pt,
    breaklines=true,
    breakatwhitespace=false,
    frame=single,
    rulecolor=\color{black},
    captionpos=b,
    keepspaces=true,
    showspaces=false,
    showstringspaces=false,
    showtabs=false,
    tabsize=2
}

\begin{document}

\bf \maketitle

\begin{center}
	\rule[0.01cm]{17cm}{0.01cm}
\end{center}

\abstract{Many low-level software have been shipped using the C programming language.
And some of them, such as EKA2 use the C++ programming language. Although notoriously difficult, one may adapt to those constraints in order to deliver one such operating system kernel.
This is why most production-grade software (Linux, XNU, and NT) are mostly written in C. With a higher-level subset in C++. However, when correctly applying the following principles to freestanding development, it becomes much easier to ensure correctness of such programs.
}

\begin{center}
	\rule[1cm]{17cm}{0.01cm}
\end{center}

\section{The Three Principles of Freestanding Development.}
\subsection{I: The Run-Time Evaluation Domain.}
{
The problem lies in the programming language runtime—which assumes an existing host. The contrary of a hosted environment is freestanding—a computing mode which doesn't expect a hosted runtime. Such programs may use the compile-time evaluation domain to achieve minimal run-time domain usage.
}

\subsection{II: The Compile-Time Evaluation Domain.}
{
One may avoid Virtual Method Tables or a runtime when possible. While focusing instead on meta-programming and compile-time features offered by C++. 
For example one may use templates to implement a scheduling policy algorithm.
One example of such implementation may be:

\begin{lstlisting}
struct FileTree final {
    static constexpr bool is_virtual_memory = false;
    static constexpr bool is_memory = false;
    static constexpr bool is_file = true;
    /// ...
};

struct MemoryTree final {
    static constexpr bool is_virtual_memory = false;
    static constexpr bool is_memory = true;
    static constexpr bool is_file = false;
    /// ...
};
\end{lstlisting} Source: \href{https://github.com/nekernel-org/nekernel/blob/develop/src/kernel/KernelKit/CoreProcessScheduler.h#L78-L105}{Link}. \\ Which is why the `constexpr' keyword is very powerful here for the Compile-Time Evaluation Domain, we avoid the many pitfalls of the Run-Time Evaluation Domain.

\subsection{III: Memory Layout and the example of C++.}
{
The Virtual Method Table (now defined as the VMT) is a big part of the problem, one may illustrate the following:

\begin{lstlisting}
/// /std:c++20 /Wall

#include <iostream>

class A
{
public:
    explicit A() = default;
    virtual ~A() = default;

    virtual void doImpl() 
    {
        std::cout << "doImpl()\r\n";
    }
};

class B : public A
{
public:
    explicit B() = default;
    ~B() override = default;
};

int main() {
    B callImpl;
    callImpl.doImpl();
}
\end{lstlisting} Source: \href{https://godbolt.org/z/aK6Y98xnd}{Link}. \\The following can instead be done to achieve similar results using the Compile-Time Evaluation Domain.

\begin{lstlisting}
inline constexpr auto kInvalidType = 0;

template <class Driver>
concept IsValidDriver = requires(Driver drv) {
	{ drv.IsActive() && drv.Type() > kInvalidType };
};
\end{lstlisting} Source: \href{https://github.com/nekernel-org/nekernel/blob/develop/src/libDDK/DriverKit/c%2B%2B/driver_base.h}{Link}.\\ Now, the problem with freestanding development is that such feature may be abused, and it is mitigated by following the TTPI.

\subsection{IV: The Three Prongs on Inheritance.}

The TTPI is a boolean framework used to evaluate whether one may consider using a Object Oriented programming language inside a freestanding program, consider the following:

\begin{itemize}
\item[1:] Is this implementable with compile-time protocols/concepts?
\item[2:] Is this implementable without three trade-off costs?
	\subitem Without violating the Runtime cost?
	\subitem The Verification cost?
	\subitem The Known-Ahead-Correctness cost?
\item[3:] Is this implementable without using a VMT?
\end{itemize}
}

\subsection{V: Compile-Time Vetting in a Freestanding Evaluation Domain.}

The following concept makes sure that the `class T' is vetted by the domain. Such properties are called `Vettable' such program in the domain makes sure that a `Container' is truly deemed fit for a Run-Time or Compile-Time Evaluation Domain. The `Vettable' structure makes use of template meta-programming in C++ to evaluate whether a `Container' shall be vetted. Such system may look as such in a Compile-Time Evaluation Domain:

\begin{lstlisting}
#define NE_VETTABLE static constexpr BOOL kVettable = YES;
#define NE_NON_VETTABLE static constexpr BOOL kVettable = NO;

template <class Type>
concept IsVettable = requires(Type) {
	(Type::kVettable);
};

/// This structure is vettable.
struct Vettable {
	NE_VETTABLE;
};

/// This structure is unvettable.
struct UnVettable {
	NE_NON_VETTABLE;
};

/// One example of a usage.
if constexpr (IsVettable<UnVettable>) {
	instVet->Vet();
} else {
	instVet->Abort();
}
\end{lstlisting} Source: \href{https://github.com/nekernel-org/nekernel/blob/develop/src/kernel/NeKit/Vettable.h}{Link}.

\section{VI: Conclusion}
{ Safe and correct development in a freestanding domain is indeed possible granted the above concepts are applied and respected.}

\section{References}

\begin{enumerate}
    \item NeKernel.org (2025). NeKernel Operating System. Available at: \url{https://nekernel.org}
    
    \item Sales, J., Tasker, M. (2005). Introducing EKA2. Symbian OS Internals. Wiley. Available at: \url{https://media.wiley.com/product_data/excerpt/47/04700252/0470025247.pdf}
    
    \item Driesen, K., Hölzle, U. (1996). The direct cost of virtual function calls in C++. \textit{OOPSLA '96}. ACM. DOI: 10.1145/236338.236369
    
    \item El Mahrouss, A. (2026). Methodology for Process and Image Computation. (v1.0.0). Zenodo. https://doi.org/10.5281/zenodo.18362503
    
    \item El Mahrouss, A. (2026). The Execution Semantics: On Axioms, Domains, and Authority. (v1.0.0). Zenodo. https://doi.org/10.5281/zenodo.18362375
\end{enumerate}

\end{document}
